% ============================================================
% Round 04: Image Editing Reward
% ============================================================

\subsection{Round 04：图像编辑 Reward}

\subsubsection{Highlight}

\begin{conceptbox}
\textbf{本章相关脉络：}
\begin{itemize}[leftmargin=1.5em]
  \item \textbf{CLIP}：arXiv 2021-02-26，用图文对比学习得到 image-text 相似度模型。论文：\url{https://arxiv.org/abs/2103.00020}，代码：\url{https://github.com/openai/CLIP}
  \item \textbf{LPIPS}：arXiv 2018-01-11，学习感知相似度，常用于图像重建 / 保真度评估。论文：\url{https://arxiv.org/abs/1801.03924}，代码：\url{https://github.com/richzhang/PerceptualSimilarity}
  \item \textbf{DINO}：arXiv 2021-04-29，自监督 ViT 表征，常用于语义保持 / identity-like feature 相似度。论文：\url{https://arxiv.org/abs/2104.14294}，代码：\url{https://github.com/facebookresearch/dino}
  \item \textbf{ImageReward}：arXiv 2023-04-12，通用 text-to-image human preference reward model。论文：\url{https://arxiv.org/abs/2304.05977}，代码：\url{https://github.com/zai-org/ImageReward}
  \item \textbf{Pick-a-Pic / PickScore}：arXiv 2023-05-02，基于真实用户偏好的 text-to-image preference 数据与 CLIP-based scorer。论文：\url{https://arxiv.org/abs/2305.01569}，代码：\url{https://github.com/yuvalkirstain/PickScore}
  \item \textbf{HPS v2}：arXiv 2023-06-15，用 HPD v2 训练 Human Preference Score v2。论文：\url{https://arxiv.org/abs/2306.09341}，代码：\url{https://github.com/tgxs002/HPSv2}
  \item \textbf{EditReward}：arXiv 2025-09-30，面向 instruction-guided image editing 的 human-aligned reward model，ICLR 2026。论文：\url{https://arxiv.org/abs/2509.26346}，代码：\url{https://github.com/TIGER-AI-Lab/EditReward}
  \item \textbf{HP-Edit / HP-Scorer}：arXiv 2026-04-21，CVPR 2026，面向图像编辑的人类偏好后训练框架。论文：\url{https://arxiv.org/abs/2604.19406}
\end{itemize}
\end{conceptbox}

\subsubsection{本章要解决的问题}

上一章 SFT 解决的是：

\[
\text{有标准目标图 } I_\text{tgt}
\quad\Rightarrow\quad
\text{让模型模仿它}
\]

但后训练真正关心的是：

\[
\text{模型生成了一个编辑结果}
\quad\Rightarrow\quad
\text{这个结果到底好不好}
\]

也就是要定义 reward：

\[
R(I_\text{src}, c, I_\text{edit})
\]

其中：

\begin{itemize}[leftmargin=1.5em]
  \item $I_\text{src}$ 是原图；
  \item $c$ 是编辑指令；
  \item $I_\text{edit}$ 是模型输出的编辑图；
  \item $R$ 是一个标量分数，越高表示越符合我们想要的编辑结果。
\end{itemize}

\begin{intubox}
\textbf{直觉：} 文本侧 reward model 给回答打分；图像编辑侧 reward model 给“原图 + 指令 + 编辑后图像”这个三元关系打分。
\end{intubox}

\subsubsection{图像编辑 Reward 的四个维度}

图像编辑不是只看好不好看。一个好的编辑结果通常要同时满足四件事：

\[
\text{指令遵循}
\quad
+
\quad
\text{未编辑区域保持}
\quad
+
\quad
\text{图像质量}
\quad
+
\quad
\text{人类偏好}
\]

所以 reward 可以抽象成：

\[
R
=
w_\text{inst}\widetilde{R}_\text{inst}
+
w_\text{pres}\widetilde{R}_\text{pres}
+
w_\text{qual}\widetilde{R}_\text{qual}
+
w_\text{pref}\widetilde{R}_\text{pref}
\]

这里加了波浪线 $\widetilde{R}$，是为了提醒：不同 reward 的量纲不一样，真实训练前通常要做 normalization / clipping / calibration。

\begin{warnbox}
\textbf{关键点：} 图像编辑 reward 不是单一指标。只优化 CLIP 可能破坏原图；只优化保持可能什么都不改；只优化美学可能把编辑任务变成重新生成。
\end{warnbox}

\subsubsection{第一类：指令遵循 Reward}

指令遵循问的是：

\[
\text{输出图像有没有完成编辑指令 } c
\]

最简单的做法是用图文相似度模型，例如 CLIP。

把输出图像编码成图像特征：

\[
h_I
=
f_I(I_\text{edit})
\]

把文本条件编码成文本特征：

\[
h_T
=
f_T(c_\text{target})
\]

然后计算 cosine similarity：

\[
R_\text{CLIP}
=
\cos(h_I,h_T)
=
\frac{
h_I^\top h_T
}{
\|h_I\|\|h_T\|
}
\]

这里我写的是 $c_\text{target}$，而不是直接写 $c$，因为编辑指令经常不是最终图像描述。

例如：

\[
c=\text{``把红车改成蓝车''}
\]

如果直接用 CLIP 比较输出图和“把红车改成蓝车”这句话，CLIP 未必能准确理解“从红到蓝”的编辑关系。更合理的目标文本可能是：

\[
c_\text{target}=\text{``a blue car''}
\]

\begin{warnbox}
\textbf{CLIP 的局限：} CLIP 更擅长判断“图像和文本是否匹配”，不擅长判断“相对于原图，是否做了正确编辑”。所以 CLIP 可以做辅助 reward，但很难单独承担图像编辑 reward。
\end{warnbox}

\subsubsection{第二类：保持 Reward}

图像编辑最容易出问题的地方是 over-edit：

\[
\text{该改的地方改了}
\quad
\text{但不该改的地方也乱变了}
\]

保持 reward 问的是：

\[
\text{未编辑区域是否还像原图}
\]

如果有 mask $m$，其中 $m=1$ 表示允许编辑区域，$1-m$ 表示应该保持区域，可以写像素级保持：

\[
R_\text{pixel-pres}
=
-
\left\|
(1-m)\odot
\left(
I_\text{edit}-I_\text{src}
\right)
\right\|_1
\]

负号表示：差异越小，reward 越高。

但像素级距离太严格。一个轻微光照变化、压缩变化、纹理变化，都会让像素差异变大。所以更常用感知距离：

\[
R_\text{LPIPS-pres}
=
-
\text{LPIPS}
\left(
(1-m)\odot I_\text{edit},
(1-m)\odot I_\text{src}
\right)
\]

如果更关心语义 / identity，而不是逐像素一致，可以用 DINO / image encoder 特征：

\[
R_\text{DINO-pres}
=
\cos
\left(
\phi(I_\text{edit}),
\phi(I_\text{src})
\right)
\]

也可以只在非编辑区域上抽特征：

\[
R_\text{region-pres}
=
\cos
\left(
\phi((1-m)\odot I_\text{edit}),
\phi((1-m)\odot I_\text{src})
\right)
\]

\begin{intubox}
\textbf{直觉：} LPIPS 更像“看起来像不像”；DINO feature 更像“语义和结构是不是还像”。它们都比直接逐像素相减更接近人的视觉感受。
\end{intubox}

\subsubsection{第三类：图像质量 Reward}

图像质量问的是：

\[
\text{输出图像是否真实、清晰、无明显伪影}
\]

可以用 aesthetic scorer、image quality scorer，或者 VLM 直接评分。

抽象写成：

\[
R_\text{qual}
=
Q(I_\text{edit})
\]

其中 $Q$ 可以判断：

\begin{itemize}[leftmargin=1.5em]
  \item 是否清晰；
  \item 是否有畸形结构；
  \item 是否有不自然边界；
  \item 光照 / 阴影 / 材质是否合理；
  \item 人脸、手、文字是否崩坏。
\end{itemize}

但质量 reward 不能单独优化。

例如指令是：

\[
\text{``把桌上的杯子移除''}
\]

模型输出一张完全不同但很漂亮的房间照片，质量很高，但编辑任务失败。

\begin{warnbox}
\textbf{质量不等于编辑正确：} $R_\text{qual}$ 只能说明图像自然，不说明它是否忠实执行了编辑指令。
\end{warnbox}

\subsubsection{第四类：Human Preference Reward}

Human preference reward 直接学习人的偏好：

\[
R_\text{pref}
=
r_\psi(I_\text{src},c,I_\text{edit})
\]

如果是 text-to-image reward model，常见输入是：

\[
r_\psi(c,I)
\]

例如 ImageReward、PickScore、HPS v2 主要面向：

\[
\text{text prompt}
\quad+\quad
\text{generated image}
\]

也就是：

\[
r_\psi(c,I_\text{gen})
\]

但图像编辑更复杂，因为必须看原图：

\[
r_\psi(I_\text{src},c,I_\text{edit})
\]

否则 reward model 不知道：

\begin{itemize}[leftmargin=1.5em]
  \item 原图里哪些内容应该保留；
  \item 指令要求改变的是哪个对象；
  \item 输出是不是只是重新生成了一张漂亮图；
  \item identity / layout / background 是否被破坏。
\end{itemize}

\begin{summarybox}
\textbf{关键差别：}\\
Text-to-image reward 看的是：
\[
(\text{prompt}, \text{generated image})
\]
Image editing reward 必须看的是：
\[
(\text{source image}, \text{instruction}, \text{edited image})
\]
\end{summarybox}

\subsubsection{Pairwise Preference：从人工偏好训练 Reward Model}

和文本 RM 一样，图像 reward model 也常从 pairwise preference 学出来。

对同一个输入：

\[
x=(I_\text{src},c)
\]

让模型生成两个候选：

\[
I_a,\quad I_b
\]

人类标注哪个更好：

\[
I_w \succ I_l
\]

其中 $I_w$ 是 winner / chosen，$I_l$ 是 loser / rejected。

Reward model 输出：

\[
r_\psi(x,I_w),
\qquad
r_\psi(x,I_l)
\]

Bradley-Terry 偏好概率是：

\[
P_\psi(I_w\succ I_l\mid x)
=
\sigma
\left(
r_\psi(x,I_w)-r_\psi(x,I_l)
\right)
\]

训练 loss 是：

\[
\mathcal{L}_\text{RM}
=
-
\log
\sigma
\left(
r_\psi(x,I_w)-r_\psi(x,I_l)
\right)
\]

这和文本侧 reward model 训练非常像，只是 $x$ 从文本 prompt 变成了：

\[
x=(I_\text{src},c)
\]

\begin{intubox}
\textbf{直觉：} Reward model 不一定要知道“绝对满分图像”是什么；它只需要学会：在同一个原图和指令下，哪个编辑结果更符合人类偏好。
\end{intubox}

\subsubsection{VLM Scorer：让视觉语言模型按 Rubric 打分}

现代图像编辑 reward 里很常见的一种做法是 VLM-as-a-judge。

输入给 VLM：

\[
\left(
I_\text{src},
I_\text{edit},
c,
\text{rubric}
\right)
\]

让它输出一个分数：

\[
s\in\{0,1,2,3,4,5\}
\]

也可以输出多维分数：

\[
\left(
s_\text{instruction},
s_\text{preservation},
s_\text{quality},
s_\text{overall}
\right)
\]

例如一个 rubric 可以写成：

\[
\begin{array}{ll}
0 &: \text{完全没有完成编辑，或图像严重损坏} \\
1 &: \text{只完成很少部分，且有明显伪影} \\
2 &: \text{部分完成，但偏离指令或破坏原图} \\
3 &: \text{大体完成，但质量或保持一般} \\
4 &: \text{较好完成，质量和保持都不错} \\
5 &: \text{完全遵循指令，保持自然，图像质量高}
\end{array}
\]

HP-Edit 里的 HP-Scorer 就是这个方向：用少量人工偏好评分数据和预训练 VLM，构造一个自动的 human preference-aligned evaluator，再把它用于偏好数据构造和 RL reward。

\begin{warnbox}
\textbf{VLM scorer 的风险：} 它比 CLIP 更懂编辑关系，但也可能有 judge bias、打分不稳定、偏好模板敏感、以及被模型学会钻空子的 reward hacking 问题。
\end{warnbox}

\subsubsection{常见 Reward / Metric 的角色}

\begin{center}
{\small
\begin{tabular}{p{2.5cm}p{3.5cm}p{4.2cm}p{3.2cm}}
\hline
\textbf{方法} & \textbf{主要看什么} & \textbf{适合用法} & \textbf{主要风险} \\
\hline
CLIP
&
图文匹配
&
粗略判断输出是否符合目标文本
&
不理解原图到编辑图的变化关系
\\
LPIPS
&
感知相似度
&
约束未编辑区域保持
&
可能惩罚合理的局部变化
\\
DINO
&
语义 / 结构特征
&
保持 identity、布局、主体语义
&
对细粒度编辑不一定敏感
\\
Aesthetic / Quality
&
美学与自然度
&
减少低质图、伪影图
&
可能鼓励重新生成漂亮图
\\
ImageReward / PickScore / HPS
&
人类偏好
&
text-to-image 偏好评价、辅助排序
&
通常不显式看原图
\\
EditReward / HP-Scorer
&
编辑结果的人类偏好
&
图像编辑偏好打分、RL reward
&
依赖标注质量、VLM judge 可靠性
\\
\hline
\end{tabular}
}
\end{center}

\subsubsection{为什么 Reward 要做 Normalization}

不同 reward 的数值范围可能完全不同：

\[
R_\text{CLIP}\in[-1,1]
\]

\[
R_\text{LPIPS}\in[0,\infty)
\]

\[
R_\text{VLM}\in\{0,1,2,3,4,5\}
\]

如果直接相加：

\[
R
=
R_\text{CLIP}
-
\text{LPIPS}
+
R_\text{VLM}
\]

某一项可能因为数值尺度大，直接支配总 reward。

所以通常要做标准化：

\[
\widetilde{R}_i
=
\frac{
R_i-\mu_i
}{
\sigma_i+\epsilon
}
\]

或者做 clipping：

\[
\widetilde{R}_i
=
\text{clip}(R_i,a_i,b_i)
\]

也可以在 GRPO 里用 group-relative advantage，让同一个 prompt / 原图下的多张候选彼此比较：

\[
A_j
=
\frac{
R_j-\text{mean}(R_1,\ldots,R_G)
}{
\text{std}(R_1,\ldots,R_G)+\epsilon
}
\]

\begin{intubox}
\textbf{直觉：} reward 的绝对值不一定可靠，但“同一个编辑任务下，哪个候选更好”通常更稳定。这也是 group-relative 方法适合图像编辑的原因之一。
\end{intubox}

\subsubsection{Reward Hacking：每个 Reward 都会被钻空子}

图像编辑 reward 的难点不是“找一个分数”，而是防止模型只优化分数、不优化真实目标。

几个典型例子：

\begin{itemize}[leftmargin=1.5em]
  \item 只优化 CLIP：模型可能把目标物体画得很明显，但破坏原图背景；
  \item 只优化 LPIPS：模型可能几乎不改图，因为不改最容易保持；
  \item 只优化 aesthetic：模型可能生成漂亮但不相关的图；
  \item 只优化 VLM scorer：模型可能学到 scorer 偏好的表面模式；
  \item 只优化 identity：模型可能拒绝执行风格、姿态、背景等合理编辑。
\end{itemize}

因此实际 reward 设计通常要做 trade-off：

\[
\text{edit strength}
\quad\leftrightarrow\quad
\text{preservation}
\]

\[
\text{instruction following}
\quad\leftrightarrow\quad
\text{image quality}
\]

\[
\text{human preference}
\quad\leftrightarrow\quad
\text{metric stability}
\]

\begin{warnbox}
\textbf{后训练的核心风险：} reward model 不是人类偏好的真相，它只是人类偏好的近似。RL 会放大 reward model 的偏差，所以 reward 设计、约束和评测必须一起看。
\end{warnbox}

\subsubsection{从 Reward 接到后面的 DPO / RL}

有了 reward 以后，后面会出现两条路线。

第一条是 preference optimization：

\[
(I_\text{src},c,I_w,I_l)
\quad\Rightarrow\quad
\text{让模型更偏向 chosen，远离 rejected}
\]

这会进入下一章：

\[
\text{Diffusion-DPO / Image Editing DPO}
\]

第二条是 RL：

\[
I_\text{edit}\sim\pi_\theta(\cdot\mid I_\text{src},c)
\]

\[
R=R(I_\text{src},c,I_\text{edit})
\]

\[
\text{用 } R \text{ 更新 diffusion / flow trajectory}
\]

这会进入后面的：

\[
\text{DDPO / Flow-GRPO / HP-Edit}
\]

\begin{summarybox}
\textbf{本章最小结论：}\\
图像编辑 reward 必须同时看原图、指令和编辑图。它通常由指令遵循、区域保持、图像质量和人类偏好组成。CLIP、LPIPS、DINO、ImageReward、PickScore、HPS、VLM scorer 都能提供部分信号，但没有任何一个单独等价于“好编辑”。后面的 DPO 和 RL，本质上都是在利用这些 reward / preference 信号继续优化编辑模型。
\end{summarybox}

\subsubsection{本轮检查题}

\begin{enumerate}[leftmargin=1.5em]
  \item 为什么图像编辑 reward 不能只看 $I_\text{edit}$？
  \item CLIP reward 在图像编辑里最大的局限是什么？
  \item LPIPS / DINO preservation reward 分别更偏向衡量什么？
  \item Text-to-image reward model 和 image editing reward model 的输入有什么区别？
  \item Pairwise preference 如何训练 reward model？
  \item 为什么多个 reward 相加前通常要 normalization？
  \item 什么是 reward hacking？图像编辑里有哪些典型例子？
\end{enumerate}
